% Current manuscript caption and methods after Q700 replacement, 2026-10-19.
    \caption{Forecast RMSE relative to the matched ECMWF ensemble mean (ENS) on a common $1.5^\circ$ grid. From left to right: 700\,hPa specific humidity (Q700), $2$\,m temperature (T2m), and 850\,hPa zonal wind (U850). Blue shows the selected GraphCast--Mamba model M1 ($1^\circ$), dashed gray its public GraphCast\_small backbone ($1^\circ$), and green the public GraphCast checkpoint ($0.25^\circ$). From 122 candidate 00\,UTC initializations every third day in 2023, each diagnostic and lead uses the identical ENS-available dates for all four forecasts (92--105 starts per point in these panels). Squared errors are pooled before taking square roots. The horizontal line marks ENS at $100\%$; lower values indicate lower RMSE. Appendix~\ref{app:common_grid} specifies the protocol.}

\paragraph{Models and forecast protocol.}
All three GraphCast-based forecasts are generated locally from fixed checkpoints. M1 is the same selected correction used in the principal evaluation, trained on 2015--2021 and selected on 2022 development data. Its frozen backbone is the public DeepMind GraphCast\_small checkpoint trained on ERA5 1979--2015, with a $1^\circ$ grid, 13 pressure levels and mesh refinement level 5. The higher-resolution reference is the public GraphCast checkpoint trained on ERA5 1979--2017, with a $0.25^\circ$ grid, 37 levels and mesh refinement level 6. Both public checkpoints include precipitation as an input and output and use the supplied GraphCast normalization statistics. Their grids contain $181\times360$ and $721\times1440$ points, respectively. Their meshes contain 10,242 and 40,962 vertices, respectively. The higher-resolution model therefore operates on much larger spatial graphs; the computational comparison concerns FLOPs and activation memory. Exact checkpoint filenames, checksums and configuration are recorded in the figure-source manifest.

Each forecast receives observed ERA5 states at $-6$ and $0$\,h, then advances through 40 six-hour leads from 6 to 240\,h using its own forecast feedback. Time-varying $0.25^\circ$ ERA5 inputs are subsampled to $1^\circ$ for GraphCast\_small and M1 and restricted to their 13 pressure levels; these models retain static fields from their existing training and evaluation store. M1 begins each forecast with zero temporal memory and feeds back its full corrected prediction. Network inference uses bfloat16, with physical fields and state boundaries in float32.

\paragraph{Dates and ENS reference.}
The candidate sample consists of 122 initializations at 00\,UTC every third day from 1 January through 30 December 2023, with verification extending through 9 January 2024. The evaluation protocol records earlier numerical-pilot exposure, including three candidate starts (15 February, 15 July and 15 November); no checkpoint selection used this comparison. ENS denotes the ECMWF IFS ensemble-mean forecast distributed by WeatherBench~2 on its conservative $121\times240$ grid. We use available 2023 forecasts from the extended 2016--2024 store.\footnote{\raggedright The ENS stores under \url{gs://weatherbench2/datasets/ifs_ens/} are \path{2016-2024-240x121_equiangular_with_poles_conservative_mean.zarr} (forecast source) and \path{2018-2022-240x121_equiangular_with_poles_conservative_mean.zarr} (layout reference).} Missing metadata and coordinate arrays were reconstructed from the complete 2018--2022 sibling store, with decoding checked against 71 overlapping forecast chunks. Missing forecast chunks remain unavailable.

At each diagnostic and lead, we retain exactly the same ENS-available initializations for all four forecasts. This gives 91--108 starts per point across the 12 diagnostics (92--105 in Figure~\ref{fig:common_grid_selected}), with 46,930 initialization--lead--diagnostic combinations per model. The matched date set can vary between diagnostics and leads; there is no single complete subset underlying every point. Selection depends only on forecast availability, not errors. Availability masks and the full candidate manifest accompany the figures.

\paragraph{Verification and metric.}
Every forecast is verified against identical ERA5 fields~\citep{hersbach2020era5} from ARCO ERA5, with source-data repairs validated against the independently curated NSF NCAR ERA5 archive. Native model predictions and $0.25^\circ$ verification fields are remapped to the common $1.5^\circ$ grid using first-order conservative spherical-cell overlap with periodic longitude; the archived ENS fields already occupy this grid. For each diagnostic and lead, we compute spatial mean squared errors with normalized cosine-latitude weights, average equally over the matched starts, and take the square root. The plotted value is $100\,\mathrm{RMSE}_{\mathrm{model}}/\mathrm{RMSE}_{\mathrm{ENS}}$, with identical dates in numerator and denominator. These physical-unit RMSE diagnostics are separate from the normalized GraphCast training loss.

\paragraph{Interpretation.}
The figures show point estimates for the specified fixed checkpoints on the available matched samples. Ten-day verification windows overlap, so the starts are not independent weather realizations. Checkpoints also differ in architecture, pressure levels and training coverage; their comparison does not isolate resolution or attribute the full advantage to memory. The matched memory controls in Section~\ref{sec:memory_ablation} address the latter question. Lower RMSE after common-grid remapping establishes neither finer-scale skill at $0.25^\circ$ nor a measured inference speedup. At day 10, M1 has lower RMSE than the public $0.25^\circ$ checkpoint on all 12 diagnostics, but higher RMSE than ENS.

\clearpage
